\newif\ifglobal

%%%%%%%%%%%%%%%%%%%%%%%
%%% Compile to view %%%
%%%%%%%%%%%%%%%%%%%%%%%

%%%%%%%%%%%%%%%%%%%%%%%%%%%%%%%%%%%%%%%%%%%%%%%%%%%
%%% Readme:                                     %%%
%%% https://www.overleaf.com/read/bwdjvfjksvjy/ %%%
%%%%%%%%%%%%%%%%%%%%%%%%%%%%%%%%%%%%%%%%%%%%%%%%%%%

% >>> M?: marks a module setting number or important notice

% >>> M4: Set {./relative-path} to external files for this module <<<
\def \modulefiles {./14-SYSTIMER}
% To add an external file, use:
% (Replace '---' with actual filename)
% '\input{\modulefiles/tables/---.tex}' for tables
% '\includegraphics{\modulefiles/figures/---}' for figures
% '\subfile{\modulefiles/---}' for register files


% >>> M1: This module is in global project? <<<
% 'YES' - uncomment; 'NO' - comment out
\globaltrue

\ifglobal
    \documentclass[main\_\_EN.tex]{subfiles}
   
\else
    \documentclass[openany, 10pt]{book}

    % >>> M2: In ./00-shared/preamble-trm-module, also find
    %         settings for visibility of todo notes, watermark <<<
    \usepackage{preamble-trm-module}

    % >>> M3: Temporary labels in English,
    %         you can add your own labels there <<<
    \myexternaldocument{./00-shared/temp-labels-en}

\fi %\ifglobal


\setcounter{secnumdepth}{4}


% >>> M5: If this module needs any updates in
% './00-shared/preamble-trm-module.sty'
% (include additional package, etc.), contact your module PM
% or create the file './00-shared/changelog.tex' make the
% necessary updates yourself and document them in this file <<<


\begin{document}


% Sets language into which pre-defined bits of text are translated
\selectlanguage{english}

% Do not delete this block
\ifglobal
% Add the following front matter if
% this module is outside of global project
\else
    \InputIfFileExists{readme}{}{}
    {\let\clearpage\relax \listoftodos}
    {\let\clearpage\relax \listofdonetodos}
    \tableofcontents
    %\listoftables
    %\listoffigures
\fi


%%%%%%%%%%%%%%%%%%%%%%%%%
%%% TRM Module Begins %%%
%%%%%%%%%%%%%%%%%%%%%%%%%

\hypertarget{systimer}{}
\chapter{System Timer (SYSTIMER)}
\label{mod:systimer}

\section{Overview}
\chipname{} provides a 52-bit timer, which can be used to generate tick interrupts for the operating system, or be used as a general timer to generate periodic interrupts or one-time interrupts. 

The timer consists of two counters: UNIT0 and UNIT1. The counter values can be monitored by three comparators COMP0, COMP1, and COMP2. See the timer block diagram on Figure \ref{fig:system-timer-structure}.

\begin{figure}[H]
    \centering
    \includegraphics[width=0.6\textwidth]{\modulefiles/figures/14-SYSTIMER-components}
    \caption{System Timer Structure}
    \label{fig:system-timer-structure}
\end{figure}

\section{Features}

The system timer has the following features:

\begin{itemize}
    \item Two 52-bit counters and three 52-bit comparators 
    \item Software accessing registers clocked by APB\_CLK
    \item CNT\_CLK used for counting, with an average frequency of 16 MHz in two counting cycles  
    \item 40 MHz XTAL\_CLK as the clock source of CNT\_CLK
    \item 52-bit alarm values (t) and 26-bit alarm periods ($\delta$t)
    \item Two modes to generate alarms:
    \begin{itemize}
        \item Target mode: only a one-time alarm is generated based on the alarm value (t)
        \item Period mode: periodic alarms are generated based on the alarm period ($\delta$t)
    \end{itemize}
    \item Three comparators generating three independent interrupts based on configured alarm value (t) or alarm period ($\delta$t)
    \item Software configuring the reference count value. For example, the system timer is able to load back the sleep time recorded by the RTC timer via software after Light-sleep
    \item Able to stall or continue running when CPU stalls or enters the on-chip-debugging mode
    \item Alarm for Event Task Matrix (ETM) event    
\end{itemize}


\section{Clock Source Selection}
The counters and comparators use XTAL\_CLK or RC\_FAST\_CLK as the clock sources. The clock source can be selected by configuring field \hyperref[fielddesc:PCRSYSTIMERFUNCCLKSEL]{PCR\_SYSTIMER\_FUNC\_CLK\_SEL} in register \hyperref[regdesc:PCRSYSTIMERFUNCCLKCONFREG]{PCR\_SYSTIMER\_FUNC\_CLK\_CONF\_REG}. After XTAL\_CLK is scaled by a fractional divider, a $f_{XTAL\_CLK}/3$ clock is generated in one count cycle and a $f_{XTAL\_CLK}/2$ clock in another count cycle. The average clock frequency is $f_{XTAL\_CLK}/2.5$, which is 16 MHz, i.e. the CNT\_CLK in Figure \ref{fig:system-timer-Alarm-Generation}. The timer counter is incremented by 1/16 $\mu s$ on each CNT\_CLK cycle. 

Software operation such as configuring registers is clocked by APB\_CLK. For more information about APB\_CLK, see Chapter \ref{mod:resclk} \textit{\nameref{mod:resclk}}. 

The following two bits of system registers are also used to control the system timer:

\begin{itemize}
    \item Set \hyperref[fielddesc:PCRSYSTIMERCLKEN]{PCR\_SYSTIMER\_CLK\_EN} in register \hyperref[regdesc:PCRSYSTIMERCONFREG]{PCR\_SYSTIMER\_CONF\_REG} to enable APB\_CLK signal to the system timer.
    \item Set \hyperref[fielddesc:PCRSYSTIMERRSTEN]{PCR\_SYSTIMER\_RST\_EN} in register \hyperref[regdesc:PCRSYSTIMERCONFREG]{PCR\_SYSTIMER\_CONF\_REG} to reset the system timer. 
\end{itemize}

Note that if the timer is reset, its registers will be restored to their default values. For more information, please refer to Chapter \ref{mod:resclk} \textit{\nameref{mod:resclk}}. 


\section{Functional Description}
\begin{figure}[H]
    \centering
    \includegraphics[width=0.8\textwidth]{\modulefiles/figures/14-SYSTIMER-architecture}
    \caption{System Timer Alarm Generation}
    \label{fig:system-timer-Alarm-Generation}
\end{figure}


Figure \ref{fig:system-timer-Alarm-Generation} shows the procedure to generate alarm in system timer. In this process, one timer counter and one timer comparator are used. An alarm interrupt will be generated accordingly based on the comparison result in comparator.


\subsection{Counter}
The system timer has two 52-bit timer counters, shown as UNIT\regindex{n} (\regindex{n} = 0 or 1). Their counting clock source is a 16 MHz clock, i.e. CNT\_CLK. Whether UNIT\regindex{n} works or not is controlled by two bits in register \hyperref[regdesc:SYSTIMERCONFREG]{SYSTIMER\_CONF\_REG}:
\begin{itemize}
    \item \hyperref[fielddesc:SYSTIMERTIMERUNIT1WORKEN]{SYSTIMER\_TIMER\_UNIT\regindex{n}\_WORK\_EN}: set this bit to enable the counter UNIT\regindex{n} in system timer.
    \item \hyperref[fielddesc:SYSTIMERTIMERUNIT1CORE0STALLEN]{SYSTIMER\_TIMER\_UNIT\regindex{n}\_CORE0\_STALL\_EN}: if this bit is set, the counter UNIT\regindex{n} stops when CPU is stalled. The counter continues its counting after the CPU resumes.
\end{itemize}
The configuration of the two bits to control the counter UNIT\regindex{n} is shown below, assuming that CPU is stalled. 

\begin{table}[H]
    \centering
    \caption{UNIT\regindex{n} Configuration Bits}
    \label{tab:systimer-config-bits}
    \begin{threeparttable}
    \begin{tabular}{ | c | c | p{5.5cm}| }
    \hline
        \rowcolor{lightgray}
        \textbf{SYSTIMER\_TIMER\_} & \textbf{SYSTIMER\_TIMER\_} & \textbf{Counter } \\
        \rowcolor{lightgray}
        \textbf{UNIT\regindex{n}\_WORK\_EN} & \textbf{UNIT\regindex{n}\_CORE0\_STALL\_EN}& \textbf{UNIT\regindex{n}} \\\hline
        % table contents
        0 & x\tnote{ *} & Not at work \\\hline
        \multirow{2}{*}{1} & \multirow{2}{*}{1} & Stop counting, but will continue its counting after the CPU resumes \\\hline
        1 & 0 & Keep counting \\\hline
        \end{tabular}

        \begin{tablenotes}
            \item[*] x: Don't-care.
        \end{tablenotes}
     \end{threeparttable}
    \end{table}

When the counter UNIT\regindex{n} is at work, the count value is incremented on each counting cycle. When the counter UNIT\regindex{n} is stopped, the count value stops increasing and keeps unchanged.

The lower 32 and higher 20 bits of initial count value are loaded from the registers \hyperref[fielddesc:SYSTIMERTIMERUNIT0LOADLO]{SYSTIMER\_TIMER\_UNIT\regindex{n}\_LOAD\_LO} and \hyperref[fielddesc:SYSTIMERTIMERUNIT0LOADHI]{SYSTIMER\_TIMER\_UNIT\regindex{n}\_LOAD\_HI}. Writing 1 to the bit \hyperref[fielddesc:SYSTIMERTIMERUNIT0LOAD]{SYSTIMER\_TIMER\_UNIT\regindex{n}\_LOAD} will trigger a reload event, and the current count value will be changed immediately. If UNIT\regindex{n} is at work, the counter will continue to count up from the new reloaded value.

Writing 1 to \hyperref[fielddesc:SYSTIMERTIMERUNIT0UPDATE]{SYSTIMER\_TIMER\_UNIT\regindex{n}\_UPDATE} will trigger an update event. The lower 32 and higher 20 bits of current count value will be locked into the registers \hyperref[fielddesc:SYSTIMERTIMERUNIT0VALUELO]{SYSTIMER\_TIMER\_UNIT\regindex{n}\_VALUE\_LO} and \hyperref[fielddesc:SYSTIMERTIMERUNIT0VALUEHI]{SYSTIMER\_TIMER\_UNIT\regindex{n}\_VALUE\_HI}, and then \hyperref[fielddesc:SYSTIMERTIMERUNIT0VALUEVALID]{SYSTIMER\_TIMER\_UNIT\regindex{n}\_VALUE\_VALID} is asserted. Before the next update event, the values of \hyperref[fielddesc:SYSTIMERTIMERUNIT0VALUELO]{SYSTIMER\_TIMER\_UNIT\regindex{n}\_VALUE\_LO} and \hyperref[fielddesc:SYSTIMERTIMERUNIT0VALUEHI]{SYSTIMER\_TIMER\_UNIT\regindex{n}\_VALUE\_HI} remain unchanged.


\subsection{Comparator and Alarm}
The system timer has three 52-bit comparators, shown as COMP\regindex{x} (\regindex{x} = 0, 1, or 2). The comparators can generate independent interrupts based on different alarm values (t) or alarm periods ($\delta$t).

Configure \hyperref[fielddesc:SYSTIMERTARGET0PERIODMODE]{SYSTIMER\_TARGET\regindex{x}\_PERIOD\_MODE} to choose from the two alarm modes for each COMP\regindex{x}:
\begin{itemize}
    \item 1: period mode
    \item 0: target mode
\end{itemize}

In period mode, the alarm period ($\delta$t) is provided by the register \hyperref[fielddesc:SYSTIMERTARGET2PERIOD]{SYSTIMER\_TARGET\regindex{x}\_PERIOD}. Assuming that current count value is t1, when it reaches (t1 + $\delta$t), an alarm interrupt will be generated. When the count value reaches (t1 + 2*$\delta$t), another alarm interrupt also will be generated. By such way, periodic alarms are generated.

In target mode, the lower 32 bits and higher 20 bits of the alarm value (t) are provided by  \hyperref[fielddesc:SYSTIMERTIMERTARGET0LO]{SYSTIMER\_TIMER\_TARGET\regindex{x}\_LO} and \hyperref[fielddesc:SYSTIMERTIMERTARGET0HI]{SYSTIMER\_TIMER\_TARGET\regindex{x}\_HI}. Assuming that current count value is t2 (t2 <= t), an alarm interrupt will be generated when the count value reaches the alarm value (t). Unlike in period mode, only one alarm interrupt is generated in target mode.

\hyperref[fielddesc:SYSTIMERTARGET0TIMERUNITSEL]{SYSTIMER\_TARGET\regindex{x}\_TIMER\_UNIT\_SEL} is used to choose the count value from which timer counter to be compared for alarm:

\begin{itemize}
        \item 1: Use the count value from UNIT\regindex{1}
        \item 0: Use the count value from UNIT\regindex{0}
\end{itemize}

Finally, set \hyperref[fielddesc:SYSTIMERTARGET0WORKEN]{SYSTIMER\_TARGET\regindex{x}\_WORK\_EN} and COMP\regindex{x} starts to compare the count value:

\begin{itemize}
        \item In target mode, COMP\regindex{x} compares with the alarm value (t).
        \item In period mode, COMP\regindex{x} compares with the alarm period (t1 + n*$\delta$t).
\end{itemize}

An alarm is generated when the count value equals to the alarm value (t) in target mode or to the start value + n*alarm period $\delta$t (n = 1, 2, 3...) in period mode. But if the alarm value (t) set in registers is less than current count value, i.e. the target has already passed, when the current count value is larger than the alarm value (t) within a range (0 \verb+~+ $2^{51}$ - 1), an alarm interrupt is also generated immediately. No matter in target mode or period mode, the low 32 bits and high 20 bits of the real alarm value can always be read from SYSTIMER\_TARGET\regindex{x}\_LO\_RO and  SYSTIMER\_TARGET\regindex{x}\_HI\_RO. The alarm trigger point and the relationship between current count value $t_{c}$ and the alarm value $t_{t}$ are shown below.

\begin{longtable}[c]{ | p{5cm} | p{7.5cm} |}
\caption{Trigger Point}\\\hline
\rowcolor{lightgray}
\textbf{Relationship Between $t_{c}$ and $t_{t}$} & \textbf{Trigger Point}\\ \hline
\endhead
$t_{c}$- $t_{t}$ <= 0 & $t_{c}$ = $t_{t}$, an alarm is triggered.\\ \hline 
0 <= $t_{c}$ - $t_{t}$ < $2^{51}$ - 1 \newline ( $t_{c}$ < $2^{51}$ and $t_{t}$ < $2^{51}$, \newline or $t_{c}$ >= $2^{51}$ and $t_{t}$ >= $2^{51}$) & An alarm is triggered immediately. \\ \hline 
\multirow{3}{*}{
$t_{c}$ - $t_{t}$ >= $2^{51}$ - 1} & 

$t_{c}$ overflows after counting to its maximum value 

52'hfffffffffffff, and then starts counting up from 0. 

When its value reaches $t_{t}$, an alarm is triggered. \\ \hline 
\end{longtable}


\subsection{Event Task Matrix}
The system timer on \chipname{} supports the Event Task Matrix (ETM) function, which allows the system timer's ETM tasks to be triggered by any peripherals' ETM events, or system timer's ETM events to trigger any peripherals' ETM tasks. This section introduces the ETM tasks and events related to the system timer. For more information, please refer to Chapter \ref{mod:evntaskmatrix} \nameref{mod:evntaskmatrix}.

The system timer can generate the following ETM event:

\begin{itemize}
\item SYSTIMER\_EVT\_CNT\_CMP\regindex{x}: Indicates the alarm pulses generated by COMP\regindex{x}.
\end{itemize}

When \hyperref[fielddesc:SYSTIMERETMEN]{SYSTIMER\_ETM\_EN} is set to 1, the alarm pulses can trigger the ETM event.

\subsection{Synchronization Operation}
The clock (APB\_CLK) used in software operation is separate from CNT\_CLK to drive the timer counters and comparators. Synchronization is needed for some configuration registers. A complete synchronization action takes two steps:

\clearpage
\begin{enumerate}
    \item Software writes specific values to configuration fields, see the first column in Table \ref{synchronization-configuration-table}.
    \item Software writes 1 to corresponding bits to start synchronization, see the second column in Table \ref{synchronization-configuration-table}.
\end{enumerate}


\begin{longtable}[c]{ | p{5.5cm} | p{6cm} |}
\caption{Synchronization Operation for Configuration Registers}\\\hline
\rowcolor{lightgray}
\label{synchronization-configuration-table}
\textbf{Configuration Fields}&\textbf{Synchronization Enable Bit}\\ \hline
\hyperref[fielddesc:SYSTIMERTIMERUNIT0LOADLO]{SYSTIMER\_TIMER\_UNIT\regindex{n}\_LOAD\_LO}

\hyperref[fielddesc:SYSTIMERTIMERUNIT0LOADHI]{SYSTIMER\_TIMER\_UNIT\regindex{n}\_LOAD\_HI} & \multirow{2}{*}{ \hyperref[fielddesc:SYSTIMERTIMERUNIT0LOAD]{SYSTIMER\_TIMER\_UNIT\regindex{n}\_LOAD}}
\\ \hline 
\hyperref[fielddesc:SYSTIMERTARGET2PERIOD]{SYSTIMER\_TARGET\regindex{x}\_PERIOD}

\hyperref[fielddesc:SYSTIMERTIMERTARGET0HI]{SYSTIMER\_TIMER\_TARGET\regindex{x}\_HI}

\hyperref[fielddesc:SYSTIMERTIMERTARGET0LO]{SYSTIMER\_TIMER\_TARGET\regindex{x}\_LO}& \multirow{3}{*}{\hyperref[fielddesc:SYSTIMERTIMERCOMP0LOAD]{SYSTIMER\_TIMER\_COMP\regindex{x}\_LOAD}} 
\\ \hline 
\end{longtable}

Synchronization is also needed for reading some status registers since the timer counter related status have a different clock than APB\_CLK.  A complete synchronization action takes three steps:
\begin{enumerate}
    \item Software writes specific values to the updating register \hyperref[fielddesc:SYSTIMERTIMERUNIT0UPDATE]{SYSTIMER\_TIMER\_UNIT\regindex{n}\_UPDATE}.
    \item Software reads the corresponding bit \hyperref[fielddesc:SYSTIMERTIMERUNIT0VALUEVALID]{SYSTIMER\_TIMER\_UNIT\regindex{n}\_VALUE\_VALID} to be valid to check synchronization is done.
    \item Software reads corresponding status registers \hyperref[fielddesc:SYSTIMERTIMERUNIT0VALUEHI]{SYSTIMER\_TIMER\_UNIT\regindex{n}\_VALUE\_HI} and \hyperref[fielddesc:SYSTIMERTIMERUNIT0VALUELO]{SYSTIMER\_TIMER\_UNIT\regindex{n}\_VALUE\_LO}.    
\end{enumerate}

\subsection{Interrupt}
Each comparator has one level-type alarm interrupt, named as SYSTIMER\_TARGET\regindex{x}\_INT. Interrupts signal is asserted high when the comparator starts to alarm. Until the interrupt is cleared by software, it remains high. To enable interrupts, set the bit \hyperref[fielddesc:SYSTIMERTARGET0INTENA]{SYSTIMER\_TARGET\regindex{x}\_INT\_ENA}.



\section{Programming Procedure}\label{sec:instructions-of-usage}
When configuring COMP\regindex{x} and UNIT\regindex{n}, please ensure the corresponding COMP and UNIT are at work.
\subsection{Read Current Count Value}\label{sec:to-read-system-timer-value}
\begin{enumerate}
    \item Set \hyperref[fielddesc:SYSTIMERTIMERUNIT0UPDATE]{SYSTIMER\_TIMER\_UNIT\regindex{n}\_UPDATE} to update the current count value of COMP\regindex{x} into \hyperref[fielddesc:SYSTIMERTIMERUNIT0VALUEHI]{SYSTIMER\_TIMER\_UNIT\regindex{n}\_VALUE\_HI} and \hyperref[fielddesc:SYSTIMERTIMERUNIT0VALUELO]{SYSTIMER\_TIMER\_UNIT\regindex{n}\_VALUE\_LO}.
    \item Poll the reading of \hyperref[fielddesc:SYSTIMERTIMERUNIT0VALUEVALID]{SYSTIMER\_TIMER\_UNIT\regindex{n}\_VALUE\_VALID} till it's 1. Then, user can read the count value from \hyperref[fielddesc:SYSTIMERTIMERUNIT0VALUEHI]{SYSTIMER\_TIMER\_UNIT\regindex{n}\_VALUE\_HI} and \hyperref[fielddesc:SYSTIMERTIMERUNIT0VALUELO]{SYSTIMER\_TIMER\_UNIT\regindex{n}\_VALUE\_LO}.
    \item Read the lower 32 bits and higher 20 bits from \hyperref[fielddesc:SYSTIMERTIMERUNIT0VALUELO]{SYSTIMER\_TIMER\_UNIT\regindex{n}\_VALUE\_LO} and \hyperref[fielddesc:SYSTIMERTIMERUNIT0VALUEHI]{SYSTIMER\_TIMER\_UNIT\regindex{n}\_VALUE\_HI} respectively. 
\end{enumerate}

\subsection{Configure One-Time Alarm in Target Mode}
\begin{enumerate}
    \item Set \hyperref[fielddesc:SYSTIMERTARGET0TIMERUNITSEL]{SYSTIMER\_TARGET\regindex{x}\_TIMER\_UNIT\_SEL} to select the counter (UNIT\regindex{0} or UNIT\regindex{1}) used for COMP\regindex{x}.
    \item Read current count value, see Section \ref{sec:to-read-system-timer-value}. This value will be used to calculate the alarm value (t) in Step 4.
    \item Clear \hyperref[fielddesc:SYSTIMERTARGET0PERIODMODE]{SYSTIMER\_TARGET\regindex{x}\_PERIOD\_MODE} to enable target mode.
    \item Set an alarm value (t), and fill its lower 32 bits to \hyperref[fielddesc:SYSTIMERTIMERTARGET0LO]{SYSTIMER\_TIMER\_TARGET\regindex{x}\_LO}, and the higher 20 bits to \hyperref[fielddesc:SYSTIMERTIMERTARGET0HI]{SYSTIMER\_TIMER\_TARGET\regindex{x}\_HI}.
    \item Set \hyperref[fielddesc:SYSTIMERTIMERCOMP0LOAD]{SYSTIMER\_TIMER\_COMP\regindex{x}\_LOAD} to synchronize the alarm value (t) to COMP\regindex{x}, i.e., load the alarm value (t) to the COMP\regindex{x}.   
    \item Set \hyperref[fielddesc:SYSTIMERTARGET0WORKEN]{SYSTIMER\_TARGET\regindex{x}\_WORK\_EN} to enable the selected COMP\regindex{x}. COMP\regindex{x} starts comparing the count value with the alarm value (t).
    \item Set \hyperref[fielddesc:SYSTIMERTARGET0INTENA]{SYSTIMER\_TARGET\regindex{x}\_INT\_ENA} to enable timer interrupt. When Unit\regindex{n} counts to the alarm value (t), a SYSTIMER\_TARGET\regindex{x}\_INT interrupt is triggered.
\end{enumerate}

\subsection{Configure Periodic Alarms in Period Mode}
\begin{enumerate}
    \item Set \hyperref[fielddesc:SYSTIMERTARGET0TIMERUNITSEL]{SYSTIMER\_TARGET\regindex{x}\_TIMER\_UNIT\_SEL} to select the counter (UNIT\regindex{0} or UNIT\regindex{1}) used for COMP\regindex{x}.
    \item Set an alarm period ($\delta$t), and fill it to \hyperref[fielddesc:SYSTIMERTARGET0PERIOD]{SYSTIMER\_TARGET\regindex{x}\_PERIOD}.
    \item Set \hyperref[fielddesc:SYSTIMERTIMERCOMP0LOAD]{SYSTIMER\_TIMER\_COMP\regindex{x}\_LOAD} to synchronize the alarm period ($\delta$t) to COMP\regindex{x}, i.e., load the alarm period ($\delta$t) to COMP\regindex{x}.
    \item Clear and then set \hyperref[fielddesc:SYSTIMERTARGET0PERIODMODE]{SYSTIMER\_TARGET\regindex{x}\_PERIOD\_MODE} to configure COMP\regindex{x} into period mode.
    \item Set \hyperref[fielddesc:SYSTIMERTARGET0WORKEN]{SYSTIMER\_TARGET\regindex{x}\_WORK\_EN} to enable the selected COMP\regindex{x}. COMP\regindex{x} starts comparing the count value with the sum of (start value + n*$\delta$t) (n = 1, 2, 3...).
    \item Set \hyperref[fielddesc:SYSTIMERTARGET0INTENA]{SYSTIMER\_TARGET\regindex{x}\_INT\_ENA} to enable timer interrupt. A SYSTIMER\_TARGET\regindex{x}\_INT interrupt is triggered when Unit\regindex{n} counts to start value + n*$\delta$t (n = 1, 2, 3...) set in Step 2.
\end{enumerate}

\subsection{Update After Light-sleep}
\begin{enumerate}
    \item Configure the RTC timer before the chip goes to Light-sleep mode, to record the exact sleep time. For more information, see Chapter \ref{mod:lowpowm} \textit{\nameref{mod:lowpowm}}.
    \item Read the sleep time from the RTC timer when the chip is woken up from Light-sleep mode. 
    \item Read current count value of system timer, see Section \ref{sec:to-read-system-timer-value}.
    \item Convert the time value recorded by the RTC timer from the clock cycles based on RTC\_SLOW\_CLK to that based on 16 MHz CNT\_CLK. For example, if the frequency of RTC\_SLOW\_CLK is 32 kHz, the recorded RTC timer value should be converted by multiplying by 500. 
    \item Add the converted RTC value to the current count value of the system timer:
    \begin{itemize}
        \item Fill the new value into \hyperref[fielddesc:SYSTIMERTIMERUNIT0LOADLO]{SYSTIMER\_TIMER\_UNIT\regindex{n}\_LOAD\_LO} (low 32 bits) and \hyperref[fielddesc:SYSTIMERTIMERUNIT0LOADHI]{SYSTIMER\_TIMER\_UNIT\regindex{n}\_LOAD\_HI} (high 20 bits).
        \item Set \hyperref[fielddesc:SYSTIMERTIMERUNIT0LOAD]{SYSTIMER\_TIMER\_UNIT\regindex{n}\_LOAD} to load new timer value into system timer. By such way, the system timer is updated.    \end{itemize}
\end{enumerate}


%%%%%%%%%%%%%%%%%%%%%%%%
%%% Register Summary %%%
%%%%%%%%%%%%%%%%%%%%%%%%

\clearpage
\hypertarget{systimer-summ}{}
\section{Register Summary}
\label{sec:28-systimer-reg-summ}

The addresses in this section are relative to system timer base address provided in Table \ref{tab:sysmem-base-address} in Chapter \ref{mod:sysmem} \textit{\nameref{mod:sysmem}}.

The abbreviations given in Column \textbf{Access} are explained in Section \hyperref[glossary-access-types]{\textit{Access Types for Registers}}.

\subfile{\modulefiles/14-SYSTIMER-reg-summ__EN}


%%%%%%%%%%%%%%%%%
%%% Registers %%%
%%%%%%%%%%%%%%%%%

\clearpage
\section{Registers}

The addresses in this section are relative to system timer base address provided in Table \ref{tab:sysmem-base-address} in Chapter \ref{mod:sysmem} \textit{\nameref{mod:sysmem}}.

\subfile{\modulefiles/14-SYSTIMER-reg__EN}


\end{document}
